% SPDX-License-Identifier: GPL-3.0-only
% Licensed under GNU GPL version 3; https://www.gnu.org/licenses/gpl-3.0.html . See LICENSE at repository root.
\section{Local generation and cyclic certification}\label{sec:certification}

\subsection{Generation produces evidence for one inference}
The public functions are
\begin{align*}
 \code{generate_axiom}&:\ J\longmapsto\mathsf{option}(\mathsf{AxiomInstance}),\\
 \code{generate_rules}&:\ J\longmapsto\mathsf{List}(\mathsf{RuleInstance}),\\
 \code{generate_requested}&:\ (J,\mathsf{request})
               \longmapsto\mathsf{option}(\mathsf{RuleInstance}).
\end{align*}
An instance contains its theorem, not just a rule name. A successful result
has exactly the requested conclusion, a well-formed source and well-formed
premises, and valid trace-edge endpoints with progress included in valid pairs.
\sources{\rocq{RuleGeneration.v}{generate_axiom},
  \rocq{RuleGeneration.v}{generate_axiom_spec},
  \rocq{RuleGeneration.v}{generate_rules_spec},
  \rocq{RuleGeneration.v}{generate_requested_spec}}

The finite scan visits left positions, then right positions, preserving their
context splits. It stops at the manual pure terminal boundary. It does not
choose arbitrary quantifier terms, discover backlinks, or synthesize a pure
proof. Explicit Boolean and quantifier requests remain usable at that boundary.
The direct primitive matches and independent matches of derived expansions
are concatenated, so overlapping abbreviations retain their named schemes.
\sources{\rocq{RuleGeneration.v}{scan_left}, \rocq{RuleGeneration.v}{scan_right},
  \rocq{RuleGeneration.v}{right_at}, \rocq{RuleGeneration.v}{left_at},
  \rocq{RuleGeneration.v}{generate_rules}}

\begin{center}
\begin{tabular}{@{}L{0.36\linewidth}L{0.58\linewidth}@{}}
\toprule
Request & Caller supplies\\
\midrule
\code{Local} & Side, zero-based position, and exact public rule name.\\
\code{Weaken}, \code{ExchangeAdjacent} & Side and position.\\
\code{Instantiate}, \code{Eigen} & Side, position, and typed term or fresh variable.\\
\code{SubstituteWorld}, \code{SubstituteTrace} & Original premise, name, and replacement.\\
\code{Cut} & Four ordered context partitions and cut occurrence.\\
\code{BacklinkExact}, \code{BacklinkTags} & Companion and, for tag changes, a finite renaming list.\\
\code{RotateLabels}, \code{UnitLabels} & Original premise, path, and direction or unit flags.\\
\code{RewriteUpdate} & Side, position, label path, state-expression depth, and law.\\
\bottomrule
\end{tabular}
\end{center}
\sources{\rocq{RuleGeneration.v}{rule_request}}

A finite tag renaming uses the first matching pair and leaves unmatched tags
unchanged. The tag-backlink request checks that renaming the source gives
the supplied companion exactly. Requests also check their constructed
conclusion against $J$, so an incompatible supplied original premise fails.
$\mathsf{None}$ means no verified instance was generated for that request;
it is not evidence that the goal is false.
\sources{\rocq{RuleGeneration.v}{rename_lookup},
  \rocq{RuleGeneration.v}{ensure_rule}, \rocq{RuleGeneration.v}{generate_requested}}

\subsection{One shared kernel instance}
Generation and certification must share the same instantiated rule kernel.
Independent applications of a Rocq functor can generate different inductive
types. A working assembly pattern is:
\begin{quote}\small
\begin{verbatim}
Module Kernel.
  Module Core := MakeRules MySignature.
  Module Model <: Core.Pr.Sem.INTERPRETATION.
    Definition interpretation := Core.Pr.Sem.Model
      Carrier chosen_inhabitants typed_functions typed_predicates.
  End Model.
  Include Core.ForModel Model.
End Kernel.
Module G := MakeRuleGeneration MySignature Kernel.
Module C := MakeCertification MySignature Kernel.
\end{verbatim}
\end{quote}
Here the four model arguments are user definitions with the types in
Section~\ref{sec:signature}. The complete instance in the heap example is a
compilable version of this assembly pattern.
\sources{\rocq{Rules.v}{KERNEL}, \rocq{Benchmarks/UserStructures.v}{Kernel}}

\subsection{Proof graphs and structural checks}
A graph stores a root judgment and an indexed finite list of nodes. Index
zero is the root. A node is open, an axiom, a rule with an ordered list of
premise targets, or a backlink with one target companion. The immutable
builder commands are \code{CloseWith}, \code{ApplyRule}, and \code{AddBacklink}.
A failed builder command leaves a sticky false success flag.
\sources{\rocq{../Cyclic/Framework.v}{node_kind},
  \rocq{../Cyclic/Framework.v}{graph}, \rocq{../Cyclic/Framework.v}{command},
  \rocq{../Cyclic/Framework.v}{execute}, \rocq{../Cyclic/Framework.v}{run_script}}

The structural check requires a nonempty graph, agreement of node zero with
the root, no open nodes, exact rule conclusion and ordered premise matches,
in-range targets, legal edge metadata, and reachability of every stored node
from zero. A backlink must have one premise, a different target index, and
tag-erased equality of its judgment with the companion. A backlink does not
receive a progress mark simply for returning to a previous judgment.
\sources{\rocq{../Cyclic/Framework.v}{node_wfb},
  \rocq{../Cyclic/Framework.v}{structural_check}}

\subsection{Relational and finite-witness checks}
The checker additionally composes finite tag relations, marking a composed
arc decreasing whenever a constituent arc decreases. For a closed graph
path, the closure check requires a decreasing diagonal in the associated
tag relation closure. This is the executable \code{relational_check}.
Its presence should not be mistaken for a proof that it accepts every
possible manuscript trace discipline.
\sources{\rocq{../Cyclic/Framework.v}{relational_check},
  \rocq{../Cyclic/Relations.v}{relation_compose},
  \rocq{../Cyclic/Relations.v}{loop_relation_ok}}

A finite witness gives each node $i$ an optional tag $t_i$ and a potential
$p_i\in\N$. A node is \emph{live} when a path with as many edges as the
graph has nodes starts there, witnessing access to a cycle in the finite
checked graph. On every edge $i\to j$ into a live target, the witness requires
\begin{center}
\begin{tabular}{@{}lll@{}}
\toprule
$t_i$ & $t_j$ & Requirement\\
\midrule
absent & absent & $p_j<p_i$\\
absent & present & No additional transition condition\\
present & absent & Rejected\\
$a$ & $b$ & $(a,b)\in V_e$ and $((a,b)\in P_e\lor p_j<p_i)$\\
\bottomrule
\end{tabular}
\end{center}
A present selected tag must belong to its live node's judgment. Edges into
nonlive regions do not need an infinite-trace witness.
\sources{\rocq{../Cyclic/Framework.v}{liveb},
  \rocq{../Cyclic/Framework.v}{trace_witness},
  \rocq{../Cyclic/Framework.v}{trace_transition},
  \rocq{../Cyclic/Framework.v}{witness_check}}

The complete executable check is
\[
 \mathsf{rawCheck}(g,tw)=
 \mathsf{structuralCheck}(g)\land\mathsf{relationalCheck}(g)
                      \land\mathsf{witnessCheck}(g,tw).
\]
\sources{\rocq{../Cyclic/Framework.v}{raw_check}}

\subsection{From Boolean acceptance to validity}
The certification bridge reflects successful structure checks into semantic
countermodel advancement. Open nodes are rejected, axioms cannot be false,
and rule/backlink nodes use their stored countermodel operations. Following
false successors for any prescribed finite length proves that every false
node is live. Witness reflection supplies the trace-transition obligations.
Together these facts construct a semantic certificate.
\sources{\rocq{Certification.v}{structurally_checked_advance},
  \rocq{Certification.v}{structurally_checked_false_live},
  \rocq{Certification.v}{witness_check_transition},
  \rocq{Certification.v}{certify_raw_check}}

\begin{theorem}[Accepted graph soundness]\label{thm:certificate}
For the fixed signature, model, and shared kernel, if
$\mathsf{rawCheck}(g,tw)=\mathsf{true}$, then $\Valid_M(\mathsf{root}(g))$.
\end{theorem}
\sources{\rocq{Certification.v}{raw_check_sound},
  \rocq{../Cyclic/Framework.v}{certificate_sound}}

\begin{proof}[Proof sketch]
Assume a false root. A false node has a false successor, and hence a live
successor. Encode the combined measure at node $i$ and world $w$ by
\[
 m(i,w)=\begin{cases}
 (1,(p_i,0)),&t_i\text{ absent},\\
 (0,(\rho_{J_i,w}(t_i),p_i)),&t_i\text{ present}.
 \end{cases}
\]
Use the strict lexicographic order on these natural-number triples. Before
tag selection, potential strictly decreases. Selecting a tag decreases the
outer component. After selection, a progressing pair strictly decreases the
semantic rank; otherwise rank cannot increase and potential strictly decreases.
The counterstep and witness therefore give a smaller measure at every false
successor, contradicting well-founded induction.
\end{proof}
\sources{\rocq{../Cyclic/Framework.v}{encode_measure},
  \rocq{../Cyclic/Framework.v}{measure_lt_wf},
  \rocq{../Cyclic/Framework.v}{certificate_step_decreases}}

The bridge obtains the semantic certificate from structural and witness
acceptance. The public check also requires the relational condition, but
the displayed soundness proof does not rely on a separately assumed
completeness or soundness theorem for relation saturation alone. A single
tag per node is sufficient for this witness format and can be conservative
relative to more general path-dependent choices.

\subsection{Trust and scope}
The development's source proofs use standard-library classical reasoning,
description/choice principles, and functional extensionality. Model fields
are explicit parameters, not hidden constructor laws. Manual leaves carry
actual Rocq proofs, and all semantic evidence in a rule or certificate is
checked by the kernel. The snapshot's validation reports record kernel
checks and assumption output; this documentation build does not rerun those
proof checks or turn the LaTeX prose into a formal proof.
